\CWHeader{Лабораторная работа \textnumero 2}

\CWProblem{
Необходимо создать программную библиотеку, реализующую указанную структуру данных, на основе которой разработать программу-словарь. В словаре каждому ключу, представляющему собой регистронезависимую последовательность букв английского алфавита длиной не более 256 символов, поставлен в соответствие некоторый номер, от 0 до $2^{64} - 1$. Разным словам может быть поставлен в соответствие один и тот же номер.

Программа должна обрабатывать строки входного файла до его окончания. Каждая строка может иметь следующий формат:

\begin{itemize}
    \item \texttt{+ word 34} --- добавить слово «word» с номером 34 в словарь. Программа должна вывести строку «OK», если операция прошла успешно, «Exist», если слово уже находится в словаре.
    \item \texttt{- word} --- удалить слово «word» из словаря. Программа должна вывести «OK», если слово существовало и было удалено, «NoSuchWord», если слово в словаре не было найдено.
    \item \texttt{word} --- найти в словаре слово «word». Программа должна вывести «OK: 34», если слово было найдено. В случае, если слово в словаре не было обнаружено, вывести «NoSuchWord».
    \item \texttt{! Save /path/to/file} --- сохранить словарь в бинарном компактном представлении на диск.
    \item \texttt{! Load /path/to/file} --- загрузить словарь из файла.
\end{itemize}

Для всех операций, в случае возникновения системной ошибки, программа должна вывести строку, начинающуюся с \texttt{ERROR:} и описывающую на английском языке возникшую ошибку.

\medskip
\textbf{Вариант структуры данных:} Красно-чёрное дерево.
}

\pagebreak

\CWHeader{Лабораторная работа \textnumero 3}

\CWProblem{
Для реализации словаря из предыдущей лабораторной работы, необходимо провести исследование скорости выполнения и потребления оперативной памяти. В случае выявления ошибок или явных недочётов, требуется их исправить.

Результатом лабораторной работы является отчёт, состоящий из:

Дневника выполнения работы, в котором отражено что и когда делалось, какие средства использовались и какие результаты были достигнуты на каждом шаге выполнения лабораторной работы.
Выводов о найденных недочётах.
Сравнение работы исправленной программы с предыдущей версией.
Общих выводов о выполнении лабораторной работы, полученном опыте.
Минимальный набор используемых средств должен содержать утилиту gprof и библиотеку dmalloc, однако их можно заменять на любые другие аналогичные или более развитые утилиты (например, Valgrind или Shark) или добавлять к ним новые (например, gcov).
}

\pagebreak